\section{Introduction}
\label{sec:intro}

How should a bottom-up model of a country's energy system be joined to a
general-equilibrium model of its economy so that each quantity crossing the boundary has
a stated economic meaning, enters the receiving model once, and can be audited?
Least-cost planning frameworks in the \CLEWS{} tradition
\citep{howells_clews_2013,howells_osemosys_2011} determine what a national energy
transition builds and operates, but they contain no households, firms, or government and
no behavioural response to the prices implied by the plan. Overlapping-generations (OLG)
general-equilibrium platforms such as \OGCore{} \citep{debacker_evans_ogcore} represent
forward-looking households, fiscal closure, and incidence across ages and lifetime-income
groups, but contain no energy technology; in the shipped model, energy is not a priced
input to production. Neither model alone can determine how a transition plan affects
output, the public budget, or the households that bear its costs.

This paper presents an auditable soft link between \CLEWS{}/\OSeMOSYS{} and \OGCore{},
explains the economic representation of each cross-model mapping, and tests the
implemented forward pass and re-solve seam in a Philippine application. We use
\emph{channel} as a compact label for a named mapping. Four channels carry outputs of a
solved \CLEWS{} scenario into \OGCore{}: the electricity-cost path, public grid
investment, generation capital intensity, and air-quality health effects. Two apply
scenario policy inputs: a carbon price and an investment tax credit. Two write
\OGCore{} outputs back to \CLEWS{}: the equilibrium return and sectoral activity. For
each mapping, the interface records the transferred quantity, the receiving-model
parameter, the accounting treatment, the admissible range, and the evidence currently
available for the mapping. The same contract supplies country-specific concordances and
units rather than embedding a particular country or industry index in the channel code.

The models remain separate programs that exchange a small set of files. This architecture
follows from their solution concepts. \OSeMOSYS{} solves one perfect-foresight linear
program over the planning horizon, whereas \OGCore{} computes an equilibrium among
overlapping cohorts along a transition path. The classic hard-link formulation, which
places a macroeconomic planner's objective inside the energy optimisation, is therefore
not available without recasting both models as a new simultaneous equilibrium system.
The implemented sequence solves the \OGCore{} baseline, solves the \CLEWS{} scenario,
applies the forward mappings, solves the economic reform, and emits the return and
activity paths. A controlled \CLEWS{} re-solve confirms that a patched demand path
survives the model boundary. The outer controller that would damp and repeat these steps
to a price--quantity fixed point is not yet implemented.

The principal economic design problem is electricity-cost transmission. The theoretically
preferred signal is the shadow price on the \OSeMOSYS{} commodity-balance constraint: the
marginal system cost of supplying another unit. In the Philippine application that dual
is degenerate, binding in scattered years and approaching zero elsewhere. The deployed
signal is therefore the reform-to-baseline ratio of a curated electricity-cost index or,
where that index is unavailable, reconstructed levelized cost of electricity. It is a
dense average-cost proxy, not a market-clearing price. This distinction limits the
interpretation of any fixed point obtained with it.

Transmitting the cost proxy also requires an explicit approximation because \OGCore{}
does not contain intermediate energy purchases. Household electricity expenditure is
represented by a revenue-neutral consumption-price wedge. Industrial electricity
expenditure is represented by an input--output-calibrated cost-push on industry
productivities. These are disjoint expenditure bases: households and firms each face the
same proportional electricity-cost change, while the social accounting matrix determines
how much of total expenditure lies on each route. Electricity self-use is removed before
the upstream cost propagation is applied. The resulting composite carries the change into
both household budgets and production costs; it does not reproduce substitution away from
energy or clear an intermediate energy market. Energy as a priced production input is the
structural extension beyond the present paper.

The OLG partner changes the questions the coupling can answer. Age-specific pollution
effects enter survival and effective labour rather than an external valuation appended to
the resource-model result. Household heterogeneity permits incidence analysis, while the
government budget makes the financing and recycling of price or investment interventions
explicit. This is a different analytical margin from neighbouring frameworks: IEEM
emphasises land use, ecosystem services, and natural-capital accounts, while CLEWs--IO
emphasises employment consequences within resource planning. The present interface
emphasises intertemporal macroeconomic adjustment, demography, fiscal closure, and
distribution across cohorts and income groups.

The Philippine application demonstrates the implemented mechanisms rather than evaluates
the underlying policy scenario. The coupled forward run combines the electricity-cost
mapping, public grid investment, and health. Output rises by about $0.2\%$ during the
first decade, consistent with the working-age morbidity improvement, and is $0.14\%$
below baseline in the long-run steady state, where higher production costs and the age
composition of the survival gain operate together. The recycled household price experiment
produces regressive incidence. Separate experiments test individual mechanisms, but they
are not an additive decomposition of the coupled run because their configurations do not
collectively reproduce it. A matched cumulative decomposition is therefore required before
the difference between the isolated and coupled results can be attributed.

\Cref{sec:background} situates the framework in the integration literature.
\Cref{sec:models} describes the parts of the two models touched by the interface.
\Cref{sec:framework} presents the price signal, architecture, interface contract, and
accounting rules. \Cref{sec:channels} describes the economic mappings and their evidence.
\Cref{sec:calibration} documents the data and calibration, and \Cref{sec:results}
reports the Philippine demonstration. \Cref{sec:discussion} states the limitations and
\Cref{sec:conclusion} the development priorities.
